%%%%%%%%%%%%%%%%%%%%%%%%%%---Abstract---%%%%%%%%%%%%%%%%%%%%%%%%%%%%%%%%%%%
\newpage
\thispagestyle{plain}
\addcontentsline{toc}{chapter}{Abstract}

\begin{center}
\huge{\textbf{Abstract}}
\end{center}
\vspace{0.3cm}

This thesis benchmarks simulated expectation-value mitigation and repetition-code
syndrome decoding. Mitigation uses circuit-grouped holdouts under idealised and
device-informed noise. On idealised random circuits, classical RBF kernel ridge
reaches $0.0404$ RMSE, fidelity kernel ridge $0.0421$, and linear zero-noise
extrapolation $0.0737$; across ten grouped holdouts the paired comparison favours RBF in nine and fidelity in none; the tenth interval includes zero. Under device-informed noise, the paired
RBF-minus-fidelity intervals on the random holdout include zero under both
noise-scaling protocols, and replacing parameter scaling by unitary folding
narrows the kernels' margin below linear extrapolation from about $0.040$ to
$0.014$ RMSE. With fidelity entries estimated from $10^5$ shots, the fidelity
kernel reaches $0.0658$ RMSE against $0.0629$ with exact entries, and each
prediction needs $3.15\times10^8$ simulated overlap executions in addition to
the $12{,}288$ input executions that both kernels share. Under phenomenological
noise, corrected matching and exact logical-class decoding improve across tested
distances at $p=0.02$. At distance three, the best valid minimum-weight matching
tie choices attain exact posterior risk at each sampled error rate.

\vspace{0.4cm}
\noindent\textbf{Keywords:} quantum error mitigation, quantum error correction,
quantum kernels, syndrome decoding.
